\cleardoublepage
\chapter*{Streszczenie}
\addcontentsline{toc}{chapter}{Streszczenie}

Rozwój procesorów wielordzeniowych stwarza możliwość przyspieszania obliczeń poprzez wykorzystanie przetwarzania równoległego. 
Samo zwiększanie liczby jednostek wykonawczych nie gwarantuje jednak proporcjonalnego wzrostu wydajności, 
ponieważ na uzyskiwane rezultaty wpływają między innymi sposób podziału danych, 
organizacja obliczeń oraz narzut związany z~tworzeniem i~synchronizacją procesów. 

Celem pracy była analiza wydajności i~efektywności wybranych algorytmów sortowania równoległego na procesorze wielordzeniowym. 
W~języku Python opracowano sekwencyjne i~równoległe implementacje algorytmów Quick Sort,
Merge Sort, Bucket Sort oraz Sample Sort z~wykorzystaniem modułu \texttt{multiprocessing} i~pamięci współdzielonej.

Badania przeprowadzono dla zbiorów od 1~000 do 1~000~000 elementów, obejmujących dane całkowite i~zmiennoprzecinkowe, 
losowe, z~duplikatami oraz częściowo posortowane.
Analizowano czas wykonania, wykorzystanie CPU i~pamięci RAM, przyspieszenie oraz efektywność. 

Dla losowego zbioru 1~000~000 liczb całkowity najkrótszy czas wykonania uzyskano dla 
Bucket Sort i~wyniósł on około 1,20~s przy 16 jednostkach wykonawczych. 
Dla Quick Sort najlepszy wynik wyniósł około 1,80~s przy 2 jednostkach, 
dla Merge Sort około 2,15~s przy 4 jednostkach, a~dla Sample Sort około 2,37~s przy 8 jednostkach. 
Najwyższe przyspieszenie uzyskano dla Sample Sort i~wyniosło około 2,60, 
podczas gdy dla Bucket Sort, Merge Sort i~Quick Sort osiągnęło odpowiednio około 1,80, 1,59 oraz 1,21.

Wyniki wykazały, że zwiększanie liczby jednostek wykonawczych nie zawsze prowadziło do skrócenia czasu wykonania, 
a~uzyskana efektywność zależała od~zastosowanego algorytmu, rozmiaru oraz charakterystyki danych wejściowych.
Największe korzyści z~przetwarzania równoległego uzyskano dla większych zbiorów danych. 
Przeprowadzone badania pozwoliły potwierdzić, że efektywne wykorzystanie 
równoległości wymaga dopasowania jej poziomu do charakterystyki algorytmu oraz przetwarzanych danych.